% Research Note of the AAS. Template: every double-brace key is filled by fill_rnaas.py
% from the result files; do not put numbers in by hand. Limits (AAS preparation
% guidelines, read 2026-10-20): 1,500 words including title, headers, captions,
% authors, acknowledgements and references, 150 of them for the abstract; one
% figure or one table; first person singular; no line numbers.
\documentclass{aastex701}

\begin{document}

\title{DEBBIE: 427,022 Galaxy Candidates with No Prior Identification in the DESI Legacy Imaging Surveys DR10}

\author{Ben Wiseman}
\affiliation{Linnet Labs PTY}
\email{benjamin.h.wiseman@gmail.com}

\begin{abstract}
I present DEBBIE, 427,022 low surface brightness galaxy candidates selected from
the DESI Legacy Imaging Surveys DR10 source catalogue, none of which the galaxy
catalogues I searched, NED, SIMBAD or DESI DR1 identify as a galaxy. Candidates are chosen
by cuts on Tractor size, surface brightness and data quality, and screened
against 1.1$\times10^{8}$ catalogued galaxies. Visual inspection of
stratified samples finds a galaxy at the catalogued position for
76--89\% of entries. Boosted trees and a bidirectional recurrent
network trained on radial light profiles score every entry with usable imaging.
No entry is removed on its score, but the 57\% with a mean score above 0.2 are
95--97\% galaxies. Within the surface-brightness window the
cuts keep 56--89\% of known galaxies with effective radii of
3--15$''$ and almost none above. All redshifts are photometric. The
designations, DEBBIE NNNNNN, are registered with the IAU.
\end{abstract}

\section{Introduction}

Imaging surveys catalogue every source they detect, and those catalogues
contain galaxies that no galaxy catalogue lists. DEBBIE selects them from the
DR10 sweep catalogues of the DESI Legacy Imaging Surveys \citep{dey2019}, which
hold a Tractor model fit \citep{lang2016} for every detection. It adds no
detections: every entry is a Tractor source, and the claim is only that none of
the catalogues and services below identifies it as a galaxy. The name
stands for DR10 Ensemble of Boosted trees and Bidirectional recurrent networks
for Identifying Extragalactic sources. A full paper is in preparation.

\section{Selection}

I keep every source in the DR10 south sweeps with \texttt{TYPE} \texttt{EXP},
\texttt{DEV}, \texttt{SER} or \texttt{REX}, which excludes point sources;
$r_{\rm eff} = \texttt{SHAPE\_R} \ge 3''$; mean surface brightness within
$r_{\rm eff}$, $\mu_{\rm eff} = g + 2.5\log_{10}(2\pi r_{\rm eff}^2)$, between
24.0 and 26.5 mag arcsec$^{-2}$, with $g$ from \texttt{FLUX\_G} and not corrected
for extinction; \texttt{NOBS\_G} and \texttt{NOBS\_R} $\ge 2$; \texttt{FLUX\_G},
\texttt{FLUX\_R} and \texttt{FLUX\_Z} $> 0$; and \texttt{MASKBITS} $= 0$. Fainter
than 26.5, Tractor increasingly fits large radii to noise. The main footprint is
$|b| \ge 30^\circ$, about 12,900 deg$^2$, and a flagged extension covers
$12^\circ < |b| < 30^\circ$ with the same cuts.

\section{Screening}

A candidate is removed if (1) a galaxy in a local index lies within
$\max(2r_{\rm eff}, 10'')$; the index holds 110,190,191 objects from HyperLEDA
\citep{makarov2014}, SDSS DR12 photometric galaxies \citep{alam2015}, GLADE+
\citep{dalya2022}, the 2MASS Extended Source Catalog \citep{jarrett2000}, 6dFGS
\citep{jones2009} and SMUDGes \citep{zaritsky2023}; (2) it lies within half the
$D_{25}$ diameter of a HyperLEDA galaxy, which removes fragments of large
galaxies; (3) NED or SIMBAD \citep{wenger2000} lists an object of a galaxy,
quasar or AGN type within $\max(5'', r_{\rm eff})$, and a failed NED query also
removes it; (4) DESI DR1 \citep{desidr1} has a spectrum with \texttt{ZWARN}
$= 0$ within $2''$, whatever its class; or (5) three or more other candidates lie
within $1.5'$, which happens on defective imaging. DES DR2 sources
\citep{abbott2021des} do not count against a candidate, since that catalogue
lists sources, not identifications. Rows within $\max(6'', r_{\rm eff})$ of each
other are merged, keeping the largest.

The catalogue has three parts that share one number sequence
(Table~\ref{tab:parts}). The original part, at $|b| \ge 30^\circ$, and
78,234 of the extension entries were screened with fixed radii in place of
rule (3), removing any candidate with a NED galaxy within $30''$, any SIMBAD
object within $15''$, or any index object, DES DR2 included, within $45''$.
The recovered part holds the candidates at $|b| \ge 30^\circ$ that those wider
radii removed and that pass rules (1)--(5); the rest of the extension was
recovered in the same way.

\section{Scores}

For each entry with usable imaging I sample the $g$, $r$ and $z$ images in 24
annuli and 32 position angles about the catalogued position and divide all three
bands by one amplitude, so that profile shape and colour survive and brightness
does not. Gradient-boosted trees and a bidirectional recurrent network score
these profiles, and the catalogue gives both scores and their mean. Trained on
the DeepShadows labels \citep{tanoglidis2021}, low surface brightness galaxies
against artefacts, the trees reach AUC 0.981 on that benchmark. Against my
visual verdicts the mean score separates galaxies from the rest with AUC
0.85 in the original part, 0.87 in the recovered part and 0.83 in
the extension. In the original part's sample, dropping entries with a mean score
below 0.06 keeps 95\% of the galaxies and removes 28\% of the rest.
A second pair, trained on HyperLEDA galaxies against stars and
faint Tractor fits, reached AUC 0.41 and 0.61 on the two samples
and is not released. No entry is removed on its score.

\section{Purity}

I judged by eye 999 objects from a sample of the original part and
2,002 from a sample of the other two parts, on Legacy Surveys colour images. Both samples are stratified over
($r_{\rm eff}$, $\mu_{\rm eff}$) cells and weighted by the entries each object
represents. A galaxy lies at the catalogued position for 89.1\% of the
original part, 76.1\% of the recovered part and 81.9\% of the
extension, leaving out objects I could not decide; Table~\ref{tab:parts} also
gives each figure with those counted as non-galaxies. Purity falls close to
catalogued galaxies, and the fixed radii kept the original part at least $45''$
from them; in the second sample it is 60\% within $20''$ of one and
84\% beyond $30''$.
Entries with a mean score above 0.2, 70\% of the original part and
48\% of the others, are 95.4--96.5\% galaxies in each part. One
person made every verdict.

\section{Completeness and caveats}

Measured on catalogued galaxies in the footprint, the cuts alone keep
56--89\% of those with $r_{\rm eff} = 3$--$15''$ and
$\mu_{\rm eff}$ in the window, but 1 of 34 above $15''$,
because Tractor splits large galaxies into fragments that fail the cuts. DEBBIE
cannot constrain large diffuse galaxies. The screen makes the catalogue uneven:
between right ascensions 150$^\circ$ and 200$^\circ$ it holds 2.1 entries
per deg$^2$ inside the SDSS footprint ($4^\circ < \delta < 20^\circ$), against
24.8 just south of it ($-10^\circ < \delta < -4^\circ$). In the original
part 63\% of the redshifts come from \citet{duncan2022} and the rest
from a five-band model I trained on those, which suits population statistics
only. The median is 0.29 and only 2,591 entries (1.6\%) lie
at $z < 0.05$, so most are distant galaxies rather than nearby diffuse dwarfs.

\section{Designations and data}

Designations take the form DEBBIE NNNNNN and are registered with the IAU
through the CDS Dictionary of Nomenclature. Numbers run from 3 to 453450
with gaps, because a number is assigned once and never reused. Positions are
Tractor positions on the Gaia frame. A flux-weighted centroid on the DR10 images
lies within 0.40$''$ of the catalogued position for half the entries, and
more than $2.5''$ away, or could not be measured, for 15\%. The
catalogue, with each entry's Tractor photometry and shape, part, scores and
centroid offset, is at Zenodo (\href{https://doi.org/10.5281/zenodo.23659669}{doi:10.5281/zenodo.23659669}).

\begin{table}[h]
\centering
\caption{The three parts. Purity: weighted share of inspected objects with a
galaxy at the catalogued position, undecided objects left out (95\% interval)
or counted as non-galaxies. AUC: mean score of the DeepShadows-trained pair
against the same verdicts. Centroid: share of entries whose centroid lies
within $1''$.\label{tab:parts}}
\begin{tabular}{lcrcccc}
\hline
Part & $|b|$ & Entries & Purity (\%) & Undecided as no (\%) & AUC & Centroid (\%) \\
\hline
Original & $\ge 30^\circ$ & 161,533 & 89.1 (86.4--91.6) & 83.7 & 0.85 & 81.9 \\
Recovered & $\ge 30^\circ$ & 141,405 & 76.1 (73.0--78.9) & 76.1 & 0.87 & 66.5 \\
Extension & 12--30$^\circ$ & 124,084 & 81.9 (79.2--84.3) & 81.8 & 0.83 & 65.5 \\
All & & 427,022 & & & & 72.0 \\
\hline
\end{tabular}
\end{table}

\begin{acknowledgments}
I thank C.~Loup of the IAU Clearing House for help with the designations. I used
data from the Legacy Surveys \citep{dey2019}, DESI DR1 \citep{desidr1}, DES DR2
\citep{abbott2021des} and SDSS-III \citep{alam2015}, whose acknowledgements are
at \url{https://www.legacysurvey.org/acknowledgment/},
\url{https://data.desi.lbl.gov/doc/acknowledgments/},
\url{https://www.darkenergysurvey.org/the-des-project/data-access/} and
\url{https://www.sdss3.org/collaboration/boiler-plate.php}. This publication
makes use of data products from the Two Micron All Sky Survey \citep{skrutskie2006},
which is a joint project of the University of Massachusetts and the Infrared
Processing and Analysis Center/California Institute of Technology, funded by the
National Aeronautics and Space Administration and the National Science
Foundation. The NASA/IPAC Extragalactic Database is funded by NASA and operated
by the California Institute of Technology. SIMBAD and VizieR are operated at
CDS, Strasbourg. I acknowledge the use of the HyperLEDA database
(\url{http://leda.univ-lyon1.fr}). I thank Clay, Raina, Lu, Trishy, Angie, Gigi, Tuna, Vesna, Joe, Ambre, and all my
other friends who have supported me in my eccentricity.
\end{acknowledgments}

\software{XGBoost, PyTorch, NumPy, SciPy, Astropy}

\begin{thebibliography}{}
\bibitem[Abbott et al.(2021)]{abbott2021des} Abbott, T. M. C., et al. 2021, ApJS, 255, 20
\bibitem[Alam et al.(2015)]{alam2015} Alam, S., et al. 2015, ApJS, 219, 12
\bibitem[D\'alya et al.(2022)]{dalya2022} D\'alya, G., et al. 2022, MNRAS, 514, 1403
\bibitem[DESI Collaboration(2026)]{desidr1} DESI Collaboration 2026, AJ, 171, 285
\bibitem[Dey et al.(2019)]{dey2019} Dey, A., et al. 2019, AJ, 157, 168
\bibitem[Duncan(2022)]{duncan2022} Duncan, K. J. 2022, MNRAS, 512, 3662
\bibitem[Jarrett et al.(2000)]{jarrett2000} Jarrett, T. H., et al. 2000, AJ, 119, 2498
\bibitem[Jones et al.(2009)]{jones2009} Jones, D. H., et al. 2009, MNRAS, 399, 683
\bibitem[Lang et al.(2016)]{lang2016} Lang, D., Hogg, D. W., \& Mykytyn, D. 2016, ascl:1604.008
\bibitem[Makarov et al.(2014)]{makarov2014} Makarov, D., et al. 2014, A\&A, 570, A13
\bibitem[Skrutskie et al.(2006)]{skrutskie2006} Skrutskie, M. F., et al. 2006, AJ, 131, 1163
\bibitem[Tanoglidis et al.(2021)]{tanoglidis2021} Tanoglidis, D., \'Ciprijanovi\'c, A., \& Drlica-Wagner, A. 2021, A\&C, 35, 100469
\bibitem[Wenger et al.(2000)]{wenger2000} Wenger, M., et al. 2000, A\&AS, 143, 9
\bibitem[Zaritsky et al.(2023)]{zaritsky2023} Zaritsky, D., et al. 2023, ApJS, 267, 27
\end{thebibliography}

\end{document}
